\documentclass{dlproblemset}
\usepackage{emoji}
\newcommand{\thetav}{\boldsymbol{\theta}}

\title{\textit{Autoencoders} e \textit{Variational Autoencoders}}
\subtitle{Lista de Exercícios}

\begin{document}

\paragraph{Instruções:} Responda às questões de múltipla escolha assinalando apenas uma alternativa. Nas questões dissertativas, apresente o raciocínio passo a passo, incluindo fórmulas e cálculos intermediários quando solicitado. Mantenha a notação dos slides: $g[\mathbf{x}, \boldsymbol{\phi}]$ para o \textit{encoder} (parâmetros $\boldsymbol{\phi}$), $f[\mathbf{z}, \thetav]$ para o \textit{decoder} (parâmetros $\thetav$), $q(\mathbf{z} \mid \mathbf{x}, \boldsymbol{\phi})$ para o posterior aproximado e $p(\mathbf{z}) = \mathcal{N}(0, \mathbf{I})$ para o \textit{prior}, conforme apresentadas em aula. Questões marcadas com \emoji{skull-and-crossbones} são consideradas difíceis.

\subsection*{Questões de Múltipla Escolha.}

\begin{enumerate}[I)]

    % ------------------------------------------------------------------
    % Q1 — Easy (AE bottleneck and trivial solutions)
    % ------------------------------------------------------------------
    \item Um pesquisador treina um \textit{autoencoder} sobre imagens $28 \times 28$ (entrada com $784$ \textit{features}) e observa erro de reconstrução próximo de zero, mas as representações latentes não capturam estrutura semântica útil. Ao inspecionar a arquitetura, ele percebe que o vetor latente possui dimensão $1024$ e que todas as camadas usam ativações lineares. Qual é a causa mais provável do problema?
    \begin{enumerate}[(a)]
        \item A função de perda escolhida não é compatível com imagens em escala de cinza.
        \item O \textit{decoder} possui parâmetros insuficientes para reconstruir a entrada.
        \item Não há gargalo de informação, então a rede aprende a copiar a entrada.
        \item A inicialização dos pesos foi feita com valores muito pequenos, causando colapso.
    \end{enumerate}
    % R: c

    % ------------------------------------------------------------------
    % Q2 — Easy (denoising AE objective)
    % ------------------------------------------------------------------
    \item Em um \textit{denoising autoencoder}, qual é o procedimento correto de treinamento por amostra?
    \begin{enumerate}[(a)]
        \item Corromper $\mathbf{x}$ gerando $\tilde{\mathbf{x}}$, passar $\tilde{\mathbf{x}}$ pela rede e comparar a saída com $\tilde{\mathbf{x}}$.
        \item Corromper $\mathbf{x}$ gerando $\tilde{\mathbf{x}}$, passar $\tilde{\mathbf{x}}$ pela rede e comparar a saída com $\mathbf{x}$ original.
        \item Passar $\mathbf{x}$ pela rede e comparar a saída com uma versão corrompida $\tilde{\mathbf{x}}$ do mesmo $\mathbf{x}$.
        \item Passar $\mathbf{x}$ pela rede e adicionar ruído na representação latente antes do \textit{decoder}.
    \end{enumerate}
    % R: b

    % ------------------------------------------------------------------
    % Q3 — Medium (AE vs VAE conceptual)
    % ------------------------------------------------------------------
    \item Considere um \textit{autoencoder} vanilla treinado para minimizar o erro de reconstrução. Por que esse modelo, sozinho, não pode ser usado para gerar amostras novas e plausíveis, ainda que reconstrua bem o conjunto de treino?
    \begin{enumerate}[(a)]
        \item O \textit{decoder} é determinístico e, portanto, não admite a injeção de ruído necessária para variabilidade.
        \item Não há distribuição conhecida sobre o latente, amostras aleatórias caem fora do suporte aprendido.
        \item O erro de reconstrução tende a zero, o que elimina qualquer diversidade na saída do \textit{decoder}.
        \item O \textit{encoder} não é diferenciável em relação a $\boldsymbol{\phi}$, impedindo a propagação do gradiente até a entrada.
    \end{enumerate}
    % R: b

    % ------------------------------------------------------------------
    % Q4 — Medium (Numerical KL for diagonal Gaussian)
    % ------------------------------------------------------------------
    \item O \textit{encoder} de um VAE com latente bidimensional ($D_z = 2$) produz, para uma entrada $\mathbf{x}$, os parâmetros $\boldsymbol{\mu} = (2,\, 0)$ e $\boldsymbol{\Sigma} = \mathrm{diag}(1,\, 1)$. Considerando $p(\mathbf{z}) = \mathcal{N}(0, \mathbf{I})$, qual é o valor de $D_{KL}\!\left[\,q(\mathbf{z} \mid \mathbf{x}, \boldsymbol{\phi})\,\|\,p(\mathbf{z})\,\right]$?
    \begin{enumerate}[(a)]
        \item $0$.
        \item $1$.
        \item $2$.
        \item $4$.
    \end{enumerate}
    % R: c

    % ------------------------------------------------------------------
    % Q5 — Medium (Reparameterization computation)
    % ------------------------------------------------------------------
    \item Durante o \textit{forward pass} de um VAE com latente bidimensional, o \textit{encoder} produz $\boldsymbol{\mu} = (1,\, -2)$ e desvios-padrão $\boldsymbol{\sigma} = (0.5,\, 1.5)$ (assumindo $\boldsymbol{\Sigma}$ diagonal). Foi amostrado $\boldsymbol{\epsilon}^{*} = (0.4,\, -0.2)$ a partir de $\mathcal{N}(0, \mathbf{I})$. Qual é o valor de $\mathbf{z}^{*}$ produzido pelo \textit{reparametrization trick}?
    \begin{enumerate}[(a)]
        \item $(1.2,\; -2.3)$.
        \item $(1.4,\; -2.2)$.
        \item $(0.9,\; -1.7)$.
        \item $(1.5,\; -1.7)$.
    \end{enumerate}
    % R: a

    % ------------------------------------------------------------------
    % Q6 — Medium (why reparameterization is needed)
    % ------------------------------------------------------------------
    \item Por que o \textit{reparametrization trick} é necessário durante o treinamento de um VAE?
    \begin{enumerate}[(a)]
        \item Para garantir que o \textit{prior} $p(\mathbf{z})$ seja exatamente $\mathcal{N}(0, \mathbf{I})$ durante toda a otimização.
        \item Para que a amostragem de $\mathbf{z}^{*}$ permaneça no caminho diferenciável em relação aos parâmetros do \textit{encoder}.
        \item Para reduzir o custo computacional da estimativa \textit{Monte Carlo} do termo de reconstrução.
        \item Para forçar a distribuição posterior $q(\mathbf{z} \mid \mathbf{x}, \boldsymbol{\phi})$ a ter média zero e variância unitária.
    \end{enumerate}
    % R: b

    % ------------------------------------------------------------------
    % Q7 — Medium (ELBO decomposition)
    % ------------------------------------------------------------------
    \item O ELBO usado no treinamento de um VAE pode ser escrito como a soma de dois termos. Qual alternativa identifica corretamente o papel de cada termo na função objetivo?
    \begin{enumerate}[(a)]
        \item Termo de reconstrução: $-D_{KL}\!\left[q(\mathbf{z}\mid \mathbf{x},\boldsymbol{\phi})\,\|\,p(\mathbf{z})\right]$. Termo de regularização: $\mathbb{E}_{q}[\log p(\mathbf{x} \mid \mathbf{z}, \thetav)]$.
        \item Termo de reconstrução: $\mathbb{E}_{q}[\log p(\mathbf{x} \mid \mathbf{z}, \thetav)]$. Termo de regularização: $-D_{KL}\!\left[q(\mathbf{z}\mid \mathbf{x},\boldsymbol{\phi})\,\|\,p(\mathbf{z})\right]$.
        \item Termo de reconstrução: $\mathbb{E}_{p(\mathbf{z})}[\log p(\mathbf{x} \mid \mathbf{z}, \thetav)]$. Termo de regularização: $D_{KL}\!\left[p(\mathbf{z})\,\|\,q(\mathbf{z}\mid \mathbf{x},\boldsymbol{\phi})\right]$.
        \item Termo de reconstrução: $\log p(\mathbf{x} \mid \thetav)$. Termo de regularização: $\mathbb{E}_{q}[\log q(\mathbf{z} \mid \mathbf{x}, \boldsymbol{\phi})]$.
    \end{enumerate}
    % R: b

    % ------------------------------------------------------------------
    % Q8 — Medium (posterior collapse identification)
    % ------------------------------------------------------------------
    \item Após treinar um VAE com \textit{decoder} muito expressivo, observa-se que $D_{KL}\!\left[q(\mathbf{z} \mid \mathbf{x}, \boldsymbol{\phi})\,\|\,p(\mathbf{z})\right]$ tende a zero para todas as entradas, enquanto o erro de reconstrução do \textit{decoder} permanece elevado. O que esse padrão indica?
    \begin{enumerate}[(a)]
        \item A taxa de aprendizado está alta demais e o gradiente do termo KL está dominando o gradiente da reconstrução.
        \item O \textit{encoder} aprendeu a igualar $q(\mathbf{z} \mid \mathbf{x}, \boldsymbol{\phi})$ ao \textit{prior} e o latente carrega pouca informação sobre $\mathbf{x}$.
        \item O ruído $\boldsymbol{\epsilon}^{*}$ do \textit{reparametrization trick} não está sendo amostrado de forma independente para cada exemplo.
        \item A dimensão do latente $D_z$ está pequena demais para representar a distribuição dos dados.
    \end{enumerate}
    % R: b

    % ------------------------------------------------------------------
    % Q9 — Medium (sampling new instances from a trained VAE)
    % ------------------------------------------------------------------
    \item Depois de treinar um VAE, deseja-se gerar uma nova amostra $\mathbf{x}^{*}$ que se assemelhe aos dados de treino. Qual procedimento corresponde à amostragem ancestral correta?
    \begin{enumerate}[(a)]
        \item Tomar um $\mathbf{x}$ qualquer do conjunto de treino, passar pelo \textit{encoder} para obter $\boldsymbol{\mu}, \boldsymbol{\Sigma}$ e produzir $\mathbf{x}^{*}$ via \textit{decoder} aplicado em $\boldsymbol{\mu}$.
        \item Amostrar $\mathbf{z}^{*} \sim \mathcal{N}(0, \mathbf{I})$ e passar $\mathbf{z}^{*}$ pelo \textit{decoder} para obter os parâmetros de $p(\mathbf{x} \mid \mathbf{z}^{*}, \thetav)$.
        \item Amostrar $\mathbf{z}^{*}$ uniformemente em um cubo de aresta $1$ centrado na origem e passar pelo \textit{decoder}.
        \item Amostrar $\mathbf{z}^{*}$ de $q(\mathbf{z} \mid \mathbf{x}, \boldsymbol{\phi})$ com $\mathbf{x}$ tirado do conjunto de treino e copiar a saída do \textit{decoder}.
    \end{enumerate}
    % R: b

    % ------------------------------------------------------------------
    % Q10 — Hard (skull) (what makes the V variational; intratable likelihood)
    % ------------------------------------------------------------------
    \item \emoji{skull-and-crossbones} Considere a fronteira entre um \textit{autoencoder} clássico e um \textit{variational autoencoder}. Qual afirmação abaixo descreve mais precisamente o papel da inferência variacional no VAE, dado que $p(\mathbf{x} \mid \thetav) = \int p(\mathbf{x} \mid \mathbf{z}, \thetav)\,p(\mathbf{z})\,d\mathbf{z}$ é, em geral, intratável?
    \begin{enumerate}[(a)]
        \item A inferência variacional substitui o termo $\log p(\mathbf{x} \mid \thetav)$ por um limite inferior tratável envolvendo $q(\mathbf{z} \mid \mathbf{x}, \boldsymbol{\phi})$, que é otimizado em conjunto com $\thetav$.
        \item A inferência variacional fornece uma forma exata de calcular $p(\mathbf{x} \mid \thetav)$ via uma única amostra de $\mathbf{z}$, dispensando integração numérica.
        \item A inferência variacional transforma o problema em uma classificação binária entre amostras reais e amostras geradas pelo \textit{decoder}.
        \item A inferência variacional garante que o \textit{encoder} aprenda exatamente a distribuição posterior verdadeira $p(\mathbf{z} \mid \mathbf{x}, \thetav)$.
    \end{enumerate}
    % R: a

    % ------------------------------------------------------------------
    % Q11 — Hard (skull) (KL sensitivity to small variance)
    % ------------------------------------------------------------------
    \item \emoji{skull-and-crossbones} Considere dois \textit{encoders} A e B de um VAE com latente bidimensional, ambos produzindo $\boldsymbol{\mu} = (0,\, 0)$. O \textit{encoder} A devolve $\boldsymbol{\Sigma}_A = \mathrm{diag}(0.01,\, 0.01)$, enquanto o B devolve $\boldsymbol{\Sigma}_B = \mathrm{diag}(1,\, 1)$. Comparando $D_{KL}\!\left[q_A \,\|\, p(\mathbf{z})\right]$ e $D_{KL}\!\left[q_B \,\|\, p(\mathbf{z})\right]$ contra $p(\mathbf{z}) = \mathcal{N}(0, \mathbf{I})$, qual afirmação está correta?
    \begin{enumerate}[(a)]
        \item $D_{KL}^{A} < D_{KL}^{B}$, pois variâncias menores aproximam $q$ de uma massa pontual e reduzem a divergência.
        \item $D_{KL}^{A} = D_{KL}^{B} = 0$, pois ambas as médias coincidem com a média do \textit{prior}.
        \item $D_{KL}^{A} > D_{KL}^{B}$, pois o termo $-\log\det\boldsymbol{\Sigma}$ cresce quando as variâncias se aproximam de zero.
        \item $D_{KL}^{A} \approx D_{KL}^{B}$, pois o termo dominante é $\boldsymbol{\mu}^{\top}\boldsymbol{\mu}$, que se anula em ambos os casos.
    \end{enumerate}
    % R: c

    % ------------------------------------------------------------------
    % Q12 — Medium (sparse vs bottleneck regularization)
    % ------------------------------------------------------------------
    \item Um \textit{autoencoder esparso} é uma variante em que a dimensão do latente pode ser \emph{maior} que a da entrada (overcomplete), mas ainda assim a rede aprende representações úteis. Como essa configuração evita a solução trivial de cópia?
    \begin{enumerate}[(a)]
        \item Penalizando a ativação do latente para que poucas unidades fiquem ativas simultaneamente.
        \item Aplicando \textit{dropout} apenas no \textit{decoder} para impedir a passagem direta da informação.
        \item Reduzindo a taxa de aprendizado sempre que o erro de reconstrução fica abaixo de um limiar.
        \item Substituindo o erro quadrático médio por entropia cruzada categórica no \textit{decoder}.
    \end{enumerate}
    % R: a

\end{enumerate}

\subsection*{Gabarito - Objetivas}
\begin{enumerate}[I)]
\item   c
\item   b
\item   b
\item   c
\item   a
\item   b
\item   b
\item   b
\item   b
\item   a
\item   c
\item   a
\end{enumerate}

\subsection*{Resoluções - Objetivas}
\color{blue}
\begin{enumerate}[I)]
    \item \textbf{(c)} Com $D_z = 1024 > 784$ e ativações lineares em todas as camadas, não existe gargalo de informação: a rede pode implementar a identidade (ou um par de mapas lineares que se cancelam) e copiar a entrada, zerando o erro sem aprender estrutura. A (a) é irrelevante para o sintoma descrito. A (b) contradiz o erro de reconstrução próximo de zero. A (d) produziria erro de reconstrução alto, o oposto do observado.

    \item \textbf{(b)} O \textit{denoising autoencoder} corrompe a entrada (gerando $\tilde{\mathbf{x}}$) mas usa o $\mathbf{x}$ limpo como alvo, forçando a rede a aprender a remover o ruído. A (a) treina a rede a reproduzir a própria entrada corrompida, sem efeito de limpeza. A (c) inverte os papéis e ensinaria a rede a adicionar ruído. A (d) descreve perturbação no latente, que não caracteriza o \textit{denoising autoencoder}.

    \item \textbf{(b)} Sem uma distribuição conhecida sobre o latente, não há como escolher $\mathbf{z}$ para alimentar o \textit{decoder}: pontos sorteados ao acaso tendem a cair fora da região efetivamente ocupada pelas codificações do treino, e o \textit{decoder} devolve saídas implausíveis. A (a) não é o impedimento, o \textit{decoder} de um VAE também é determinístico dado $\mathbf{z}$. A (c) é um \textit{non sequitur}, reconstruir bem não elimina diversidade. A (d) parte de premissa falsa, o \textit{encoder} é diferenciável.

    \item \textbf{(c)} $D_{KL} = \tfrac{1}{2}\left(\mathrm{Tr}[\boldsymbol{\Sigma}] + \boldsymbol{\mu}^{\top}\boldsymbol{\mu} - D_z - \log\det\boldsymbol{\Sigma}\right) = \tfrac{1}{2}(2 + 4 - 2 - 0) = 2$. A (a) corresponde a ignorar a média (com $\boldsymbol{\mu} = \mathbf{0}$ e $\boldsymbol{\Sigma} = \mathbf{I}$ a divergência seria nula). A (b) corresponde a usar $\|\boldsymbol{\mu}\| = 2$ sem elevar ao quadrado. A (d) corresponde a usar $\boldsymbol{\mu}^{\top}\boldsymbol{\mu} = 4$ sozinho, sem o fator $\tfrac{1}{2}$ e sem os demais termos.

    \item \textbf{(a)} $\mathbf{z}^{*} = \boldsymbol{\mu} + \boldsymbol{\sigma} \odot \boldsymbol{\epsilon}^{*} = (1 + 0.5 \cdot 0.4,\; -2 + 1.5 \cdot (-0.2)) = (1.2,\; -2.3)$. A (b) soma $\boldsymbol{\epsilon}^{*}$ diretamente, sem escalar por $\sigma$. As alternativas (c) e (d) resultam de trocas de componentes ou de sinais na aplicação elemento a elemento.

    \item \textbf{(b)} O nó de amostragem é estocástico e não admite derivada em relação a $\boldsymbol{\phi}$. Reescrevendo $\mathbf{z}^{*} = \boldsymbol{\mu} + \boldsymbol{\Sigma}^{1/2}\boldsymbol{\epsilon}^{*}$, a aleatoriedade fica isolada em $\boldsymbol{\epsilon}^{*}$ (amostrado de uma distribuição fixa) e o caminho por $\boldsymbol{\mu}$ e $\boldsymbol{\Sigma}^{1/2}$ permanece diferenciável. A (a) confunde o \textit{prior}, fixado por hipótese, com algo imposto pelo truque. A (c) não procede, o custo da estimativa não muda. A (d) descreveria um colapso forçado do posterior, que não é o objetivo.

    \item \textbf{(b)} O ELBO é $\mathbb{E}_{q}[\log p(\mathbf{x} \mid \mathbf{z}, \thetav)] - D_{KL}\!\left[q(\mathbf{z} \mid \mathbf{x}, \boldsymbol{\phi})\,\|\,p(\mathbf{z})\right]$: o primeiro termo mede a qualidade da reconstrução, o segundo regulariza $q$ na direção do \textit{prior}. A (a) troca os papéis dos dois termos. A (c) usa a esperança sob o \textit{prior} e inverte os argumentos da KL. A (d) confunde o ELBO com a própria verossimilhança marginal, que é intratável.

    \item \textbf{(b)} KL próxima de zero com reconstrução ruim é sinal do \textit{posterior collapse}: $q(\mathbf{z} \mid \mathbf{x}, \boldsymbol{\phi}) \approx p(\mathbf{z})$ para todo $\mathbf{x}$ e o latente deixa de informar o \textit{decoder}. A (a) descreveria instabilidade geral de otimização, não esse padrão específico. A (c) não produz esse sintoma. A (d) implicaria KL tipicamente alta, não nula.

    \item \textbf{(b)} A amostragem ancestral segue o modelo gerador: $\mathbf{z}^{*} \sim p(\mathbf{z}) = \mathcal{N}(0, \mathbf{I})$ e, em seguida, o \textit{decoder} produz os parâmetros de $p(\mathbf{x} \mid \mathbf{z}^{*}, \thetav)$. A (a) descreve reconstrução de um dado existente, não geração. A (c) amostra da distribuição errada. A (d) condiciona em um exemplo do treino, tendendo a reproduzi-lo em vez de gerar algo novo.

    \item \textbf{(a)} Como $p(\mathbf{x} \mid \thetav)$ envolve uma integral intratável, a inferência variacional introduz $q(\mathbf{z} \mid \mathbf{x}, \boldsymbol{\phi})$ e otimiza o ELBO, um limite inferior tratável de $\log p(\mathbf{x} \mid \thetav)$, conjuntamente em $\boldsymbol{\phi}$ e $\thetav$. A (b) superestima o que uma única amostra de $\mathbf{z}$ fornece. A (c) descreve uma GAN. A (d) é falsa, $q$ apenas aproxima o posterior verdadeiro.

    \item \textbf{(c)} $D_{KL}^{A} = \tfrac{1}{2}\left(0.02 + 0 - 2 - \ln(10^{-4})\right) = \tfrac{1}{2}(0.02 - 2 + 9.21) \approx 3.6$, enquanto $D_{KL}^{B} = \tfrac{1}{2}(2 + 0 - 2 - 0) = 0$. O termo $-\log\det\boldsymbol{\Sigma}$ explode quando as variâncias tendem a zero, logo $D_{KL}^{A} > D_{KL}^{B}$. A (a) inverte a intuição, uma massa quase pontual está longe de uma gaussiana de variância unitária. A (b) considera apenas as médias e ignora as covariâncias. A (d) ignora o termo de determinante, que é justamente o dominante aqui.

    \item \textbf{(a)} A penalidade de esparsidade força poucas unidades ativas por instância, criando um gargalo de capacidade efetiva mesmo com latente \textit{overcomplete}: a rede precisa selecionar o que representar. A (b) não impede a cópia pelo restante das conexões. A (c) e a (d) alteram a dinâmica de otimização e a função de perda, mas nenhuma das duas remove a solução de identidade.
\end{enumerate}
\color{black}

\newpage

\subsection*{Questões Dissertativas.}

\begin{enumerate}[I)]

    \item Considere um \textit{autoencoder} aplicado a imagens em escala de cinza com $D_x = 784$ \textit{pixels} (entrada vetorizada de $28 \times 28$). Para cada uma das configurações abaixo, indique se o modelo é \emph{undercomplete} ou \emph{overcomplete}, e discuta brevemente o efeito esperado sobre a representação aprendida.
    \begin{enumerate}[(a)]
        \item Latente com $D_z = 32$ unidades, sem outras regularizações.
        \item Latente com $D_z = 2000$ unidades, sem outras regularizações.
        \item Latente com $D_z = 2000$ unidades, treinado com penalidade de esparsidade na ativação do latente.
    \end{enumerate}

    \item Seja um VAE com latente $D_z = 2$. Para uma entrada $\mathbf{x}$, o \textit{encoder} produz $\boldsymbol{\mu} = (1,\, -1)$ e $\boldsymbol{\Sigma} = \mathrm{diag}(0.5,\, 2)$. Calcule, passo a passo, $D_{KL}\!\left[\,q(\mathbf{z} \mid \mathbf{x}, \boldsymbol{\phi})\,\|\,\mathcal{N}(0, \mathbf{I})\,\right]$ usando a fórmula fechada apresentada em aula. Em seguida, comente quais dos quatro termos ($\mathrm{Tr}[\boldsymbol{\Sigma}]$, $\boldsymbol{\mu}^{\top}\boldsymbol{\mu}$, $D_z$, $\log\det\boldsymbol{\Sigma}$) mais contribuíram para o valor obtido.

    \item Explique o \textit{reparametrization trick} respondendo aos itens abaixo.
    \begin{enumerate}[(a)]
        \item Qual é o problema do ponto de vista da \textit{backpropagation} ao amostrar $\mathbf{z}^{*} \sim q(\mathbf{z} \mid \mathbf{x}, \boldsymbol{\phi})$ diretamente?
        \item Escreva a forma reparametrizada de $\mathbf{z}^{*}$ usada em VAEs com posterior gaussiano de média $\boldsymbol{\mu}$ e covariância diagonal $\boldsymbol{\Sigma}$, identificando claramente quais variáveis ficam no caminho diferenciável e qual variável é a fonte de estocasticidade.
        \item \emoji{skull-and-crossbones} Mostre, usando a forma reparametrizada, que $\mathbb{E}[\mathbf{z}^{*}] = \boldsymbol{\mu}$ e que $\mathrm{Cov}[\mathbf{z}^{*}] = \boldsymbol{\Sigma}$.
    \end{enumerate}

    \item Compare um \textit{autoencoder} clássico e um \textit{variational autoencoder} respondendo:
    \begin{enumerate}[(a)]
        \item O que cada \textit{encoder} produz para uma mesma entrada $\mathbf{x}$?
        \item Por que somente o VAE permite \emph{amostragem ancestral} de novos $\mathbf{x}^{*}$ após o treinamento?
        \item Qual é a função objetivo otimizada em cada caso e por que o VAE não otimiza diretamente $\log p(\mathbf{x} \mid \thetav)$?
    \end{enumerate}

    \item Considere um VAE treinado em MNIST cuja loss de treino apresenta o seguinte comportamento ao longo das épocas: o termo $D_{KL}\!\left[q(\mathbf{z} \mid \mathbf{x}, \boldsymbol{\phi})\,\|\,p(\mathbf{z})\right]$ cai rapidamente para perto de zero, enquanto o termo de reconstrução estagna em um valor alto. As amostras geradas a partir de $\mathbf{z}^{*} \sim p(\mathbf{z})$ resultam em imagens de aparência média, com pouca variação entre amostras distintas.
    \begin{enumerate}[(a)]
        \item Nomeie o fenômeno observado e explique o que ele significa em termos da informação que o latente $\mathbf{z}$ carrega sobre $\mathbf{x}$.
        \item Por que esse fenômeno é especialmente comum quando o \textit{decoder} é muito expressivo (por exemplo, autoregressivo)?
        \item Cite uma estratégia conhecida para mitigar o problema e explique, em uma frase, como ela age sobre o ELBO.
    \end{enumerate}

\end{enumerate}

\subsection*{Gabarito}
\color{blue}
\begin{enumerate}[I)]

    \item Configurações.
    \begin{enumerate}[(a)]
        \item \emph{Undercomplete} ($D_z = 32 < 784$). O gargalo força o \textit{encoder} a comprimir, descartando informações redundantes; tende a produzir representações úteis para tarefas posteriores.
        \item \emph{Overcomplete} ($D_z = 2000 > 784$), sem regularização. A rede pode aprender uma cópia trivial ou indexar as instâncias do treino (memorização), com representação latente pouco informativa.
        \item \emph{Overcomplete} ($D_z = 2000$), porém regularizado por esparsidade. A penalidade obriga que apenas poucas unidades do latente fiquem ativas por instância, recuperando a utilidade da representação mesmo sem gargalo dimensional.
    \end{enumerate}

    \item Cálculo passo a passo.
    \begin{align*}
        \mathrm{Tr}[\boldsymbol{\Sigma}]   &= 0.5 + 2 = 2.5 \\
        \boldsymbol{\mu}^{\top}\boldsymbol{\mu}         &= 1^{2} + (-1)^{2} = 2 \\
        \log\det\boldsymbol{\Sigma}        &= \log(0.5) + \log(2) \approx -0.693 + 0.693 = 0 \\
        D_{z}                 &= 2 \\
        D_{KL} &= \tfrac{1}{2}\bigl(\mathrm{Tr}[\boldsymbol{\Sigma}] + \boldsymbol{\mu}^{\top}\boldsymbol{\mu} - D_z - \log\det\boldsymbol{\Sigma}\bigr) \\
               &= \tfrac{1}{2}\bigl(2.5 + 2 - 2 - 0\bigr) = \tfrac{1}{2}\,(2.5) = 1.25.
    \end{align*}
    Os dois termos que mais contribuem são $\mathrm{Tr}[\boldsymbol{\Sigma}] = 2.5$ (dispersão acima do \textit{prior} unitário, sobretudo pela segunda coordenada com $\sigma^{2} = 2$) e $\boldsymbol{\mu}^{\top}\boldsymbol{\mu} = 2$ (desvio da origem). O termo $\log\det\boldsymbol{\Sigma}$ acaba se anulando porque as variâncias $0.5$ e $2$ se compensam em escala logarítmica.

    \item \textit{Reparametrization trick}.
    \begin{enumerate}[(a)]
        \item A operação de amostragem é estocástica e não tem derivada bem definida com respeito aos parâmetros que produzem a distribuição. Sem reparametrização, o gradiente de $\mathcal{L}$ não consegue fluir do \textit{decoder} de volta até $\boldsymbol{\phi}$ através do nó de amostragem.
        \item Escrevemos $\mathbf{z}^{*} = \boldsymbol{\mu} + \boldsymbol{\Sigma}^{1/2}\,\boldsymbol{\epsilon}^{*}$, com $\boldsymbol{\epsilon}^{*} \sim \mathcal{N}(0, \mathbf{I})$. As quantidades $\boldsymbol{\mu}$ e $\boldsymbol{\Sigma}^{1/2}$ (saídas do \textit{encoder}, função de $\boldsymbol{\phi}$) permanecem no caminho diferenciável; a estocasticidade fica concentrada em $\boldsymbol{\epsilon}^{*}$, que é amostrado de uma distribuição fixa e \emph{não} depende de $\boldsymbol{\phi}$.
        \item Como $\boldsymbol{\epsilon}^{*} \sim \mathcal{N}(0, \mathbf{I})$ é independente de $\boldsymbol{\mu}$ e $\boldsymbol{\Sigma}$:
        \begin{align*}
            \mathbb{E}[\mathbf{z}^{*}]  &= \boldsymbol{\mu} + \boldsymbol{\Sigma}^{1/2}\,\mathbb{E}[\boldsymbol{\epsilon}^{*}] = \boldsymbol{\mu} + \mathbf{0} = \boldsymbol{\mu}. \\
            \mathrm{Cov}[\mathbf{z}^{*}] &= \mathbb{E}\bigl[(\mathbf{z}^{*} - \boldsymbol{\mu})(\mathbf{z}^{*} - \boldsymbol{\mu})^{\top}\bigr] \\
                                 &= \mathbb{E}\bigl[\boldsymbol{\Sigma}^{1/2}\,\boldsymbol{\epsilon}^{*}\boldsymbol{\epsilon}^{*\,\top}\,\boldsymbol{\Sigma}^{1/2}\bigr] \\
                                 &= \boldsymbol{\Sigma}^{1/2}\,\mathbb{E}[\boldsymbol{\epsilon}^{*}\boldsymbol{\epsilon}^{*\,\top}]\,\boldsymbol{\Sigma}^{1/2} \\
                                 &= \boldsymbol{\Sigma}^{1/2}\,\mathbf{I}\,\boldsymbol{\Sigma}^{1/2} = \boldsymbol{\Sigma}.
        \end{align*}
        Logo $\mathbf{z}^{*}$ é distribuído como $\mathcal{N}(\boldsymbol{\mu}, \boldsymbol{\Sigma})$, conforme desejado.
    \end{enumerate}

    \item Comparação AE versus VAE.
    \begin{enumerate}[(a)]
        \item O \textit{encoder} de um AE clássico produz um \emph{ponto} no espaço latente, $\mathbf{z} = g[\mathbf{x}, \boldsymbol{\phi}]$, determinístico dada a entrada. O \textit{encoder} de um VAE produz os parâmetros de uma \emph{distribuição} sobre o latente, $q(\mathbf{z} \mid \mathbf{x}, \boldsymbol{\phi}) = \mathcal{N}(\boldsymbol{\mu}, \boldsymbol{\Sigma})$.
        \item No AE clássico não há uma distribuição conhecida sobre o latente, então amostrar $\mathbf{z}$ ao acaso pode cair em regiões sem suporte aprendido pelo \textit{decoder}. No VAE o termo KL aproxima $q(\mathbf{z} \mid \mathbf{x}, \boldsymbol{\phi})$ do \textit{prior} $p(\mathbf{z})$, garantindo que amostras $\mathbf{z}^{*} \sim p(\mathbf{z})$ permaneçam em regiões plausíveis para o \textit{decoder}.
        \item O AE clássico minimiza apenas o erro de reconstrução, $J(\mathbf{x}, f[g[\mathbf{x}, \boldsymbol{\phi}], \thetav])$. O VAE maximiza o ELBO, que é um limite inferior de $\log p(\mathbf{x} \mid \thetav)$. Não otimizamos $\log p(\mathbf{x} \mid \thetav)$ diretamente porque a marginal $p(\mathbf{x} \mid \thetav) = \int p(\mathbf{x} \mid \mathbf{z}, \thetav)\,p(\mathbf{z})\,d\mathbf{z}$ envolve uma integral intratável devido à não linearidade do \textit{decoder}.
    \end{enumerate}

    \item Diagnóstico do treino.
    \begin{enumerate}[(a)]
        \item O fenômeno é o \emph{posterior collapse}: o \textit{encoder} aprendeu a igualar $q(\mathbf{z} \mid \mathbf{x}, \boldsymbol{\phi})$ ao \textit{prior} para qualquer $\mathbf{x}$, de modo que o latente $\mathbf{z}$ deixa de carregar informação sobre a entrada. A consequência é que amostras geradas a partir do \textit{prior} se concentram em torno de uma reconstrução média, com pouca diversidade.
        \item Quando o \textit{decoder} é poderoso o bastante para modelar $p(\mathbf{x} \mid \mathbf{z}, \thetav)$ ignorando $\mathbf{z}$ (por exemplo, modelos autoregressivos sobre os \textit{pixels}), o otimizador encontra um ótimo barato em zerar o termo KL: a contribuição de $\mathbf{z}$ para o ELBO fica marginal e o gradiente empurra $q$ na direção de $p(\mathbf{z})$.
        \item Estratégias comuns incluem \textit{KL annealing} (introduzir o termo KL gradualmente, começando com peso baixo) e \textit{free bits} (não penalizar a KL abaixo de um limiar por dimensão do latente). Ambas reduzem a pressão do termo de regularização nas fases iniciais, permitindo que o \textit{decoder} aprenda a usar $\mathbf{z}$ antes que o KL force o colapso.
    \end{enumerate}

\end{enumerate}
\color{black}

\end{document}
